\documentclass[11pt,a4paper]{article}

\usepackage{amsmath,amssymb,amsthm}
\usepackage{mathtools}
\usepackage{geometry}
\usepackage{hyperref}
\usepackage{algorithm}
\usepackage{algpseudocode}

\geometry{margin=2.5cm}

\newtheorem{definition}{Definition}
\newtheorem{theorem}{Theorem}
\newtheorem{lemma}{Lemma}
\newtheorem{proposition}{Proposition}
\newtheorem{corollary}{Corollary}
\newtheorem{remark}{Remark}
\newtheorem{assumption}{Assumption}

\title{Mathematical Foundations of LIFA-Fuzz:\\A Neural-Mathematical Protocol Fuzzer}
\author{}
\date{}

\begin{document}
\maketitle

\section{Overview}

LIFA-Fuzz is a protocol fuzzer that combines statistical signal processing
with large language model (LLM) inference to infer and exploit the grammar of
network protocols. This document presents the mathematical foundations of the
system, independent of implementation:

\begin{enumerate}
  \item \textbf{Cross-packet differential analysis} --- a suite of four
        statistical tests (Shannon entropy, population variance, Pearson
        correlation, Kendall's tau-b) that jointly classify each byte offset
        in a packet stream into a semantic field type.
  \item \textbf{EWMA adaptive sampling} --- an exponentially weighted moving
        average controller that dynamically adjusts the system's response
        sampling interval in response to measured coverage discovery rates.
  \item \textbf{Probabilistic protocol state machine inference} --- a
        black-box, specification-free method that infers a protocol's state
        machine from captured network traces using statistical message-unit
        extraction, PAM clustering, and DFA construction (Veritas-inspired).
  \item \textbf{Inverse-frequency power schedule (IFPS)} --- an
        acceptance-rejection sampling scheme for seed selection that
        allocates mutation effort inversely to past selection frequency,
        ensuring $O(1)$ expected time per selection regardless of corpus
        size.
\end{enumerate}

\section{Foundation 1: Cross-Packet Differential Analysis}

The system infers protocol structure by applying four statistical tests to
each byte offset across a collection of captured packets. Each test captures
a different structural property: randomness (entropy), dispersion (variance),
length-dependence (correlation), and sequential trend (rank correlation).

\subsection{Data Model}

Let $\mathcal{P} = \{p_1, \dots, p_N\}$ be a multiset of packets, each
$p_j \in \{0,1\}^{m_j}$ a byte string of length $m_j$. For byte offset $i$,
define the observation vector

\[
V_i = \bigl[\,p_j[i] \mid j \in \{1,\dots,N\} \land i < m_j\,\bigr],
\qquad n_i = |V_i|,
\]

where each $v \in V_i$ lies in $\{0,\dots,255\}$.

\subsection{Shannon Entropy}

\begin{definition}[Empirical entropy]
\[
H(V_i) = -\sum_{v=0}^{255} \hat p(v) \log_2 \hat p(v),\qquad
\hat p(v) = \frac{|\{x \in V_i : x = v\}|}{n_i}.
\]
\end{definition}

\begin{theorem}[Structural discrimination]
For any $V_i$ with $n_i \geq 1$,
\[
0 \leq H(V_i) \leq 8.
\]
The lower bound $H = 0$ occurs when all bytes are identical (constant
fields: magic bytes, padding). The upper bound $H = 8$ occurs when bytes
are uniformly distributed (encrypted or compressed payload). Entropy thus
separates three structural regimes:

\[
\ell_H(V_i) =
\begin{cases}
\text{STATIC}        & H \leq 0.1, \\[2pt]
\text{LOW\_ENTROPY}  & 0.1 < H < 3.5, \\[2pt]
\text{HIGH\_ENTROPY} & H \geq 3.5.
\end{cases}
\]
\end{theorem}

\begin{proof}
Nonnegativity is immediate from the definition. The upper bound follows
from Gibbs' inequality: for any distribution on a finite set $\mathcal{X}$,
$H(p) \leq \log_2 |\mathcal{X}|$, with equality iff $p$ is uniform.
Here $|\mathcal{X}| = 256$, so $\log_2 256 = 8$.
\end{proof}

\subsection{Population Variance}

\begin{definition}
\[
\sigma^2(V_i) = \frac{1}{n_i} \sum_{v \in V_i} (v - \mu_i)^2,\qquad
\mu_i = \frac{1}{n_i} \sum_{v \in V_i} v.
\]
\end{definition}

\begin{proposition}
For byte-valued data,
\[
\max \sigma^2 = \frac{255^2}{4} = 16256.25,
\]
achieved when $V_i$ contains only $0$ and $255$ in equal proportion.
\end{proposition}

Variance is used as a tiebreaker when entropy is intermediate: high-variance
low-entropy offsets suggest flag fields or enums with extreme values.

\subsection{Pearson Correlation with Packet Length}

Length fields are detected by correlating decoded byte values with packet
length. For each offset $i$, decode $V_i$ under byte widths
$k \in \{1,2,4\}$ and both endiannesses $e \in \{\text{le}, \text{be}\}$,
producing vectors $D_{i,k,e}$. Let $L_j = m_j$ be the length of packet $j$.

\begin{definition}[Length correlation]
\[
r_{i,k,e} = \frac{\sum_j (d_j - \bar d)(L_j - \bar L)}
                 {\sqrt{\sum_j (d_j - \bar d)^2}\;
                  \sqrt{\sum_j (L_j - \bar L)^2}},
\qquad
r_L(V_i) = \max_{k,e} |r_{i,k,e}|.
\]
The maximum is subject to the plausibility constraint
$\max_j D_{i,k,e}^{(j)} \leq (\max_j L_j - h) \cdot 9/8$,
where $h$ is a candidate header size.
\end{definition}

\begin{theorem}[Range]
For any finite non-degenerate sequences, $r_{i,k,e} \in [-1, 1]$.
\end{theorem}

\begin{proof}
Cauchy-Schwarz inequality applied to the centred vectors.
\end{proof}

If $r_L(V_i) \geq 0.85$, the offset is classified as CALCULATED with
sub-type LENGTH\_FIELD.

\subsection{Kendall's Tau-b for Monotonic Trend}

Sequence numbers are detected by measuring monotonic trend robust to ties.
Let $(v_j, j)$ be (value, capture-index) pairs at offset $i$.

\begin{definition}[Kendall tau-b]
For all $\binom{n_i}{2}$ unordered pairs $(j_1, j_2)$ with $j_1 < j_2$,
define:
\begin{align*}
C   &= \#\{\text{concordant}: v_{j_1} < v_{j_2}\}, \\
D   &= \#\{\text{discordant}: v_{j_1} > v_{j_2}\}, \\
T_y &= \#\{\text{tied in value}: v_{j_1} = v_{j_2}\}.
\end{align*}
Then
\[
\tau_b(V_i) = \frac{C - D}{\sqrt{(n_0 - T_y)\, n_0}},\qquad
n_0 = \binom{n_i}{2}.
\]
\end{definition}

\begin{theorem}[Computational complexity]
$D$ can be computed in $O(n_i \log n_i)$ time via merge-sort inversion
count. If $\pi$ stably sorts $V_i$ by value, then $D$ equals the number of
inversions in $\pi$ (pairs $(\alpha, \beta)$ with $\alpha < \beta$ but
$\pi(\alpha) > \pi(\beta)$). Each tie group of size $k$ contributes
$\binom{k}{2}$ to $T_y$.
\end{theorem}

If $\tau_b(V_i) \geq 0.75$, the offset is classified as CALCULATED with
sub-type SEQUENCE\_NUMBER.

\subsection{Combined Classification}

The four statistics are combined via strict priority:

\[
\ell(V_i) =
\begin{cases}
\text{STATIC}                & H \leq 0.1, \\
\text{CALCULATED (LENGTH)}  & \text{else if } r_L \geq 0.85, \\
\text{CALCULATED (SEQUENCE)} & \text{else if } \tau_b \geq 0.75, \\
\text{HIGH\_ENTROPY}         & \text{else if } H \geq 3.5, \\
\text{LOW\_ENTROPY}          & \text{otherwise}.
\end{cases}
\]

Adjacent offsets with the same label are grouped into field groups.
Multi-byte CALCULATED groups with combined width $\in \{2,3,4,8\}$ are
merged (matching uint16/uint24/uint32/uint64).

\section{Foundation 2: EWMA Adaptive Sampling}

The system adjusts its response sampling interval $k(t)$ (number of packets
between successive \texttt{recv()} calls) using an Exponentially Weighted
Moving Average (EWMA) of the coverage discovery rate.

\subsection{Proxy Metric}

Let epoch $t$ have duration $\Delta t$ seconds. The coverage increment
combines two proxy measurements:

\[
\Delta C_t = \frac{0.3 \cdot \Delta A_t + 0.7 \cdot B_t}{\Delta t},
\]

where $\Delta A_t$ is the number of newly discovered field groups and $B_t$
the number of unique hex response prefixes observed during epoch $t$.

\subsection{EWMA Recurrence and Sampling Interval}

\begin{definition}[Smoothed coverage intensity]
\[
\lambda_C(0) = 0,\qquad
\lambda_C(t) = \delta \cdot \Delta C_t + (1 - \delta) \cdot \lambda_C(t-1),
\quad \delta = 0.1.
\]
\end{definition}

\begin{definition}[Adaptive sampling interval]
\[
k(t) = \max\!\Bigl(k_{\min},\;
       \bigl\lfloor \tfrac{K_{\max}}{1 + \theta \cdot \lambda_C(t)}
       \bigr\rfloor \Bigr),
\qquad
K_{\max} = 200,\; \theta = 2.0,\; k_{\min} = 5.
\]
\end{definition}

When $\lambda_C$ is high (many recent discoveries), $k$ decreases, increasing
sampling frequency. When the protocol is saturated ($\lambda_C \to 0$),
$k$ approaches $K_{\max}$, conserving LLM budget.

\begin{theorem}[Convergence under stationarity]
If $\Delta C_t$ is stationary with mean $\mu$, then
$\mathbb{E}[\lambda_C(t)] \to \mu$ as $t \to \infty$, with effective
window $1/\delta = 10$ epochs.
\end{theorem}

\begin{proof}
Expanding: $\lambda_C(t) = \delta \sum_{i=0}^{t-1} (1-\delta)^i \Delta C_{t-i}
+ (1-\delta)^t \lambda_C(0)$. Taking expectations,
$\mathbb{E}[\lambda_C(t)] \to \mu \cdot \delta \cdot 1/\delta = \mu$.
\end{proof}

\subsection{Boundedness and Monotonicity}

The EWMA controller is a simple adaptive filter, not a feedback regulator
with stability concerns. Its properties follow directly from the formula:

\begin{theorem}[Boundedness]
$0 \leq \lambda_C(t) \leq \max_t \Delta C_t$ for all $t$.
\end{theorem}

\begin{proof}
$\lambda_C(0) = 0$. Each update is a convex combination:
$\lambda_C(t) = \delta \Delta C_t + (1-\delta)\lambda_C(t-1)$.
Since $\Delta C_t \geq 0$ and $\delta \in (0,1)$, the result follows by
induction.
\end{proof}

\begin{theorem}[Monotonic response]
$k(t)$ is a strictly decreasing function of $\lambda_C(t)$:
$\frac{\partial k}{\partial \lambda_C} = \frac{-\theta K_{\max}}
{(1+\theta \lambda_C)^2} < 0$.
\end{theorem}

\begin{proof}
Direct differentiation of $k = \lfloor K_{\max}/(1+\theta\lambda_C) \rfloor$.
\end{proof}

The controller is thus free of oscillation or chattering by construction
($k$ is a smooth, bounded, monotonic function of a bounded input), unlike
step-function AIMD controllers that jump between $k=1$ and $k=K_{\max}$.

\subsection{Regime Classification}

The normalised rate $\rho(t) = k(t) / K_{\max}$ defines the operating regime:

\[
\text{regime}(\rho) =
\begin{cases}
\text{INTENSIVE} & \rho < 0.1, \\[2pt]
\text{ACTIVE}    & 0.1 \leq \rho < \frac13, \\[2pt]
\text{NORMAL}    & \frac13 \leq \rho < \frac23, \\[2pt]
\text{SPARSE}    & \rho \geq \frac23.
\end{cases}
\]

\section{Foundation 3: Probabilistic Protocol State Machine Inference}

For stateful protocols (FTP, SSH, TLS), a fuzzer must reach deep protocol
states to exercise vulnerable code. LIFA-Fuzz infers a Probabilistic
Protocol State Machine (P-PSM) from captured traffic using a purely
statistical, specification-free approach inspired by Veritas
\cite{veritas}. No source code, RFC, or protocol keyword is required.

\subsection{Message-Unit Extraction}

Let $\mathcal{P} = \{p_1, \dots, p_N\}$ be captured packets. Define a
\emph{message unit} as a length-$\ell$ byte subsequence starting at offset
$i$ in a packet header (first $n$ bytes):

\[
u_{j,i} = p_j[i\,..\,i+\ell-1], \qquad \ell = 3,\; n = 12.
\]

\begin{definition}[Frequency-filtered unit set]
Partition $\mathcal{P}$ into two equal halves $\mathcal{P}_A, \mathcal{P}_B$.
Count the frequency $f_u^A, f_u^B$ of each distinct unit $u$ in each half.
The candidate unit set is
\[
\Phi = \{u : f_u^A \geq \lambda \;\land\; f_u^B \geq \lambda\},
\]
where $\lambda$ is the smallest threshold such that the two frequency
distributions pass a two-sample Kolmogorov--Smirnov test at significance
level $1 - \alpha < 10^{-8}$.
\end{definition}

This K-S filter retains only units whose frequency is \emph{consistently
high} across independent samples --- a signature of protocol structure
rather than random payload coincidence.

\subsection{Format Message Reconstruction}

For each packet $p_j$, greedily extend from offset $0$ using candidate
units from $\Phi$:

\[
m_j = p_j[0\,..\,i^*], \quad i^* = \max\{i : \forall\, 0 \leq j < i-\ell+1,\;
p_j[j\,..\,j+\ell-1] \in \Phi\}.
\]

The resulting $m_j$ (if $|m_j| \geq \ell$) is the packet's \emph{format
message} --- the protocol-keyword prefix that identifies its type.

\subsection{State Message Clustering}

\begin{definition}[Byte-frequency feature vector]
For a format message $m$, define
\[
\mathbf{v}(m) = \bigl[\,c_0, c_1, \dots, c_{255}\,\bigr], \qquad
c_b = |\{i : m[i] = b\}|,
\]
the 256-dimensional histogram of byte values.
\end{definition}

\begin{definition}[Jaccard distance]
\[
d(m_a, m_b) = 1 - \frac{|\{b : c_b(m_a) > 0\} \cap \{b : c_b(m_b) > 0\}|}
                    {|\{b : c_b(m_a) > 0\} \cup \{b : c_b(m_b) > 0\}|}.
\]
\end{definition}

Format messages are grouped via Partitioning Around Medoids (PAM) under
Jaccard distance. The optimal cluster count $k$ maximises the Dunn index:

\begin{definition}[Dunn index]
\[
D(k) = \frac{\min_{i \neq j}\;\delta(C_i, C_j)}
            {\max_{i}\;\Delta(C_i)},
\]
where $\delta(C_i, C_j) = d(\text{medoid}_i, \text{medoid}_j)$ is the
inter-cluster distance and $\Delta(C_i) = \max_{a,b \in C_i} d(a, b)$ is
the intra-cluster diameter.
\end{definition}

Each cluster's medoid (the member minimising total distance to all others)
is a \emph{protocol state message} --- the representative header pattern for
that state type.

\subsection{State Machine Construction}

Each captured packet is labelled with its nearest medoid state type. Flows
(session sequences) become state-type sequences $F = (\pi_1, \dots,
\pi_h)$. Transition probabilities are estimated empirically:

\begin{definition}[P-PSM transition probability]
\[
P(\pi_i \to \pi_j) = \frac{|\{t : \pi_t = \pi_i \;\land\; \pi_{t+1} =
\pi_j\}|}{|\{t : \pi_t = \pi_i\}|}.
\]
\end{definition}

The resulting directed graph $(Q, \Sigma, \delta, P)$ is a Probabilistic
Protocol State Machine: states $Q$ are the cluster medoids, alphabet
$\Sigma$ is the set of observed state-message types, transitions $\delta$
are the observed edges, and $P$ assigns probabilities.

\begin{theorem}[State labelling cost]
Labelling a new packet requires computing $k$ Jaccard distances (one per
medoid), each $O(256)$. Total: $O(256k)$ per packet, where $k$ is the
number of inferred states (typically $5$--$15$). This is fast enough for
real-time use in the Fast Loop.
\end{theorem}

\begin{proof}
Each Jaccard distance computes set intersection/union over at most 256
elements. With $k$ medoids, the total is $256k$ operations, each $O(1)$
amortised.
\end{proof}

\subsection{Black-Box Property}

The entire pipeline --- unit extraction, K-S filtering, PAM clustering, DFA
construction --- operates on raw captured bytes without any protocol
specification, source code, keyword list, or port number. It is applicable
to any stateful protocol whose traffic can be captured, making it a genuine
black-box state inference method.

\section{Foundation 4: Inverse-Frequency Power Schedule (IFPS)}

Seed selection uses an acceptance-rejection scheme that allocates mutation
effort inversely to past selection frequency, ensuring fair exploration of
the corpus.

\subsection{Energy Function}

Let $\mathcal{C} = \{s_1, \dots, s_N\}$ be the seed corpus and $f_i$ the
number of times seed $s_i$ has been selected.

\begin{definition}[Energy]
\[
E(i) = \frac{1}{f_i + 1},\qquad
E'(i) =
\begin{cases}
5 \cdot E(i) & \text{if seed $i$ is STATE\_NOVELTY}, \\
E(i)         & \text{otherwise}.
\end{cases}
\]
\end{definition}

A seed never selected ($f_i = 0$) receives maximal energy $E = 1$.
Seeds that discovered a novel state transition receive a $5\times$ boost.
As $f_i \to \infty$, $E \to 0$, ensuring underrepresented sequences are
explored before over-fuzzed ones.

\subsection{Acceptance-Rejection Sampler}

\begin{algorithm}[H]
\caption{IFPS seed selection}
\begin{algorithmic}[1]
\State $n \gets |\mathcal{C}|$
\If{$n = 0$} \Return error \EndIf
\If{$n = 1$} \Return $\mathcal{C}[0]$ \EndIf
\For{$\text{attempt} = 1 \textbf{ to } n + 10$}
    \State $i \sim \text{Uniform}\{0, \dots, n-1\}$
    \State $u \sim \text{Uniform}[0, 1]$
    \If{$u < \min(E'(i), 1)$}
        \State \Return $\mathcal{C}[i]$
    \EndIf
\EndFor
\State \Return $\mathcal{C}[\text{Uniform}\{0, \dots, n-1\}]$
\end{algorithmic}
\end{algorithm}

\begin{theorem}[Correctness]
Assume the fallback is never reached (probability $\to 0$ as $n$ grows).
Then the probability that seed $i$ is selected in any given trial is
proportional to $E'(i)$:

\[
P(\text{select } i) = \frac{E'(i)}{\sum_{j=1}^n E'(j)}.
\]
\end{theorem}

\begin{proof}
In each trial, pick $i$ uniformly: $P(\text{pick } i) = 1/n$. Accept with
probability $a_i = \min(E'(i), 1)$. By the acceptance-rejection lemma,
conditioned on acceptance occurring,

\[
P(\text{return } i) = \frac{(1/n) \cdot a_i}{\sum_j (1/n) \cdot a_j}
= \frac{a_i}{\sum_j a_j}.
\]

When $E'(i) \leq 1$ (all seeds with $f_i \geq 4$ for novel seeds or
$f_i \geq 0$ for normal seeds), $a_i = E'(i)$, so $P(\text{return } i)
\propto E'(i)$. Novel seeds with $f_i < 4$ have $E'(i) > 1$, so $a_i = 1$;
this caps acceptance at 1, ensuring novel seeds are always accepted and
equally likely among themselves, preserving proportionality in practice.
\end{proof}

\subsection{Expected Time Analysis}

Let $a_{\min} = \min_j a_j$. Since $E'(j) > 0$ for all finite $f_j$,
$a_{\min} > 0$.

\begin{theorem}[$O(1)$ expected time]
The expected number of trials per selection is at most $1 / a_{\min}$, a
constant independent of corpus size $n$. In the worst case (all seeds at the
frequency cap $f_j = 999$), $a_j = 1/1000$, giving at most $1000$ expected
trials. The probability of reaching the $n+10$ fallback is
$(1 - 1/1000)^{n+10} \to 0$ exponentially as $n$ grows.
\end{theorem}

\subsection{Corpus Maintenance}

To prevent unbounded memory growth:

\[
|\mathcal{C}| > 5000 \;\Longrightarrow\; \text{evict } 20\% \text{ most-frequent seeds},
\qquad
f_i \gets \min(f_i, 999) \text{ every } 10^4 \text{ sends}.
\]

The cap $f_i \leq 999$ prevents any single seed from dominating
indefinitely while providing a uniform lower bound $a_j \geq 1/1000$ for
the acceptance probability, guaranteeing the $O(1)$ expected-time bound.

\end{document}
